\hebrewsection{Machine Learning in Signal Processing}

The integration of Machine Learning (ML) and Deep Learning (DL) with traditional DSP is creating a new paradigm for "learned" signal processing.

\subsection{Learned Physical Layers}
Instead of hand-designing modulators and demodulators, we can use an \textbf{Autoencoder} to learn the optimal end-to-end communication system. The encoder acts as the transmitter and the decoder as the receiver, optimized over a differentiable channel model.

\subsection{Signal Classification with CNNs}
Convolutional Neural Networks (CNNs) are highly effective at classifying signals represented in the time-frequency domain (spectrograms).
\textit{Applications:}
\begin{itemize}
    \item \textbf{Modulation Recognition:} Identifying the modulation type of an intercepted signal.
    \item \textbf{Signal Denoising:} Removing non-stationary noise that traditional filters cannot handle.
\end{itemize}

\subsection{Deep Learning for Beamforming}
In massive MIMO systems, calculating the optimal beamforming weights is computationally expensive. Deep Neural Networks (DNNs) can be trained to approximate these calculations in real-time, significantly reducing latency.

\subsection{Generative Models for Channel Modeling}
Generative Adversarial Networks (GANs) can be trained to produce realistic synthetic channel responses that match the statistics of specific environments (e.g., a factory floor or a moving train), improving the robustness of system simulations.

\subsection{Challenges and Future Directions}
\begin{itemize}
    \item \textbf{Interpretability:} Deep models are often "black boxes" compared to mathematically rigorous DSP.
    \item \textbf{Complexity:} Running large neural networks on embedded DSPs or FPGAs requires specialized hardware acceleration.
    \item \textbf{Data Requirements:} Training requires massive amounts of labeled signal data.
\end{itemize}

\begin{pythonbox}{Simple CNN for Signal Classification}
\begin{verbatim}
import tensorflow as tf
from tensorflow.keras import layers

model = tf.keras.Sequential([
    layers.Conv2D(32, (3, 3), activation='relu', input_shape=(128, 128, 1)),
    layers.MaxPooling2D((2, 2)),
    layers.Flatten(),
    layers.Dense(64, activation='relu'),
    layers.Dense(10, activation='softmax')
])
\end{verbatim}
\end{pythonbox}
